\section{The Plan Grows a Parallel Node}
\label{sec:ch12-the-plan-grows-a-parallel-node}
\index{query plan!at six subgraphs|(}%

Six schemas composed on the first attempt, with no edit to any of them. That is
worth one sentence and not more: the three settings in
section~\ref{sec:ch12-what-a-service-costs-to-make} are what makes it true, and
chapter~\ref{ch:composition} already spent a chapter on what happens when they
are not.

\index{composition!errors are a fact about a schema set}%
One consequence of the graph growing does need saying, because it makes an
earlier chapter's listings unreproducible at this tag.
Chapter~\ref{ch:composition}'s error catalogue was produced by breaking a
subgraph called \code{mosaic}, which no longer exists, so the same six mistakes
are now made in \code{pricing} and the composer says different words about three
of them. The shareability errors name four subgraphs where they used to name
one, and the unsatisfiable-key error names two unresolvable fields where it used
to name four. Those listings reproduce at tag \code{ch11} and earlier, which is
what a chapter tag is for. The lesson is smaller than it looks and worth having
anyway: a catalogue of composition errors is a fact about a particular set of
subgraphs rather than a fact about the composer, and it goes stale when the set
changes.

What is worth more than a sentence is the plan. Here is the storefront query
with reviews and their authors in it, sent with
\code{X-WG-Include-Query-Plan} and \code{X-WG-Skip-Loader}, flattened to the
shape rather than printed in full:

\begin{minted}[fontsize=\footnotesize,breaklines]{text}
Sequence
  Single catalog id=0 depends=[] path=
  Parallel
    BatchEntity pricing   id=1 depends=[0] path=browseProducts.nodes
    BatchEntity inventory id=2 depends=[0] path=browseProducts.nodes
    BatchEntity reviews   id=3 depends=[0] path=browseProducts.nodes
  BatchEntity accounts    id=4 depends=[3] path=browseProducts.nodes.@.reviews.nodes.@.author
\end{minted}

Two fetches became five, in three steps, and
figure~\ref{fig:ch12-query-plan} is the same thing as a picture.

\begin{figure}[htbp]
  \centering
  % The storefront query plan at tag ch11 and at tag ch12. Referenced from
% chapter 12. Bare tikzpicture; the figure environment, caption and label live
% at the call site.
%
% The point of the picture is the middle node on the right. Three of the four
% new fetches happen at once because they depend on the same key and on nothing
% else, and the fourth cannot, because its keys do not exist until the fetch
% above it has answered. Depth in the graph is sequencing in the plan.
\begin{tikzpicture}[
    font=\small,
    head/.style={font=\small\bfseries, align=center},
    fetch/.style={draw, rounded corners=2pt, minimum width=20mm,
                  minimum height=5.5mm, align=center, font=\scriptsize},
    par/.style={draw, dashed, rounded corners=2pt, gray},
    flow/.style={-{Stealth[length=2mm]}},
    note/.style={font=\scriptsize, gray, align=center, inner sep=1pt}
  ]

  % -- left: two fetches ------------------------------------------------------
  \node[head] at (1.5,0.8) {two, at \code{ch11}};

  \node[fetch] (a1) at (1.5,0.0)  {\code{catalog}};
  \node[fetch] (a2) at (1.5,-1.5) {\code{mosaic}};

  \draw[flow] (a1) -- (a2);
  \node[note, anchor=west] at (2.7,-0.75) {one key};

  % -- right: five fetches in three steps -------------------------------------
  \node[head] at (8.4,0.8) {five, at \code{ch12}};

  \node[fetch] (b1) at (8.4,0.0)   {\code{catalog}};

  \draw[par] (5.6,-1.9) rectangle (11.2,-1.1);
  \node[fetch] (bp) at (6.7,-1.5)  {\code{pricing}};
  \node[fetch] (bi) at (8.4,-1.5)  {\code{inventory}};
  \node[fetch] (br) at (10.1,-1.5) {\code{reviews}};

  \node[fetch] (ba) at (10.1,-3.1) {\code{accounts}};

  \draw[flow] (b1) -- (bp);
  \draw[flow] (b1) -- (bi);
  \draw[flow] (b1) -- (br);
  \draw[flow] (br) -- (ba);

  \node[note, anchor=east, align=right] at (5.5,-1.5)
    {at once: same\\key, nothing else};
  \node[note, anchor=east, align=right] at (8.9,-3.1)
    {after: these keys do not\\exist until reviews answers};

\end{tikzpicture}

  \caption{The storefront plan before and after. The middle node on the right is
  the one worth having: three subgraphs depend on the same key and on nothing
  else, so the planner runs them together. The fourth cannot join them, because
  its keys are customer identifiers that do not exist until Reviews has
  answered. Depth in the graph is sequencing in the plan.}
  \label{fig:ch12-query-plan}
\end{figure}

\index{query plan!Parallel node@\code{Parallel} node}%
The \code{Parallel} node is the first one in this book, and it is the clearest
evidence that the cut was made in the right places. Three services contribute
fields to \code{Product}; all three need the same key and nothing else; so all
three go at once. A graph where the planner had to sequence those would be a
graph where the domains had been carved along the wrong seam, and the plan would
say so before any measurement did.

\index{query plan!path@\code{path}!two list traversals}%
The fourth fetch is the interesting one. Its path is the longest in this book:

\begin{minted}[fontsize=\footnotesize,breaklines]{text}
browseProducts.nodes.@.reviews.nodes.@.author
\end{minted}

Two \code{@} steps, which is chapter~\ref{ch:entity-resolution-done-right}'s
notation for a walk through a list, twice over. And it depends on fetch 3 rather
than on fetch 0, because the keys it sends are customer identifiers and nothing
knows what those are until Reviews has answered. That single dependency is the
one place Mosaic's domain graph has any depth at all, and the plan found it
without being told.

Every dependency in the plan is a key. The router collection walks all of them
and asserts \code{isKey} and a field name of \code{id} on every one, which is
the return on chapter~\ref{ch:hotchocolate-quickly}'s rule that no entity holds
a navigation property to a type another domain owns.

\subsection*{What it cost}

\index{measurement!the storefront across six services}%
The storefront query at twenty-five products, warm, through the router, read off
each subgraph's own request timeline. Counts rather than timings, which is why
they are in the gate at all: chapter~\ref{ch:enter-the-router} drew that line
and chapter~\ref{ch:entity-resolution-done-right} kept it.

This is not the query the plan above was drawn from. It is
chapter~\ref{ch:entity-resolution-done-right}'s measurement query, which stops
short of \code{Product.reviews} on purpose, because a resolver runs per review
author and the Postman collection submits a review earlier in the same gate: a
count that walked the reviews would be true only of a database nothing had
written to. Four services answer it rather than five, and the same query at the
previous tag is what makes the two columns comparable at all.

\begin{minted}[fontsize=\footnotesize,breaklines]{graphql}
{ browseProducts(first: 25) { nodes { title price { amount currency }
    shippingCost { amount currency } availableQuantity averageRating } } }
\end{minted}

\begin{minted}[fontsize=\footnotesize]{text}
                without shippingCost      with shippingCost
                resolvers      SQL        resolvers      SQL
catalog                 1        1                1        1
pricing                26        1               51        1
inventory              26        1               26        1
reviews                26        1               26        1
accounts                -        -                -        -
ordering                -        -                -        -
total                  79        4              104        4
\end{minted}

Chapter~\ref{ch:entity-resolution-done-right} printed 76 and 101 resolvers and 3
statements for the same query. The arithmetic between the two tags is where the
chapter's real claim is, and it has three parts.

\index{measurement!what an extraction moves and what it adds}%
101 became 103 across the three contributing services, because there are three
\code{\_entities} root fields where there was one. Two extra resolver
invocations for two extra HTTP round trips: that is what the split added, and it
is a rounding error against the work itself.

Catalog's row is the one chapter~\ref{ch:entity-resolution-done-right} could not
print. One resolver and one statement, invisible at that tag because Catalog had
no instrumentation, which chapter~\ref{ch:enter-the-router} called half of every
federated query going unmeasured. It was not half. It was one statement out of
four, and the interesting part is that nobody could have told you which.

\textbf{So the statement count did not move at all.} Four after, and four
before, although the two fours are not the same kind of number and it is worth
being exact about that. Three of them were measured at
tag~\code{ch11}. The fourth was not, because Catalog could not count then. What
lets me claim it anyway is that the two queries in \code{CatalogService} are
unchanged between the tags and the only edit to that file was the counter
itself, so the statement being counted here is the statement that was running
there. An extraction moves work between
processes and this one added no database work whatsoever, which is the thing I
would want to know before agreeing to do it.

The two silent rows are asserted too. Accounts and Ordering answer nothing for
this query and the gate checks that they answer nothing, which is the assertion
that would catch a router fetching from a service the query never mentions.

\subsection*{The number that moved for a strange reason}

\index{resolver count!measures tasks rather than fields}%
One number fell off a cliff, and chasing it turned up something about the count
itself that has been true since chapter~\ref{ch:the-life-of-a-request} and that
I had never noticed.

The \code{\_entities} call the gate sends Reviews directly, twenty-five products
and every review they have:

\begin{minted}[fontsize=\footnotesize]{text}
                     at ch11        at ch12
resolvers               146             26
SQL                       2              1
lookups                   2              1
\end{minted}

The statement is the author batch leaving for Accounts, which is what an
extraction does and needs no explanation. The resolvers do.

The engine still produces a hundred and twenty authors. What it no longer does
is run a hundred and twenty resolver \emph{tasks} to get them.
\code{ReviewNode.GetAuthor} takes a parent and returns a new object with no
await in it, so HotChocolate compiles it to a \code{PureFieldDelegate} and runs
it inline. Read at tag \code{16.6.0}, commit \code{8fea46e}, the compiler makes
that choice in one condition:

\begin{minted}[fontsize=\footnotesize,breaklines]{csharp}
if (field.PureResolver is not null && selection.Directives.Count == 0)
{
    return field.PureResolver;
}
\end{minted}

And the event the count comes from is raised in exactly two places, both of them
inside a task:

\begin{minted}[fontsize=\scriptsize,breaklines]{text}
src/HotChocolate/Core/src/Types/Execution/Processing/Tasks/ResolverTask.Execute.cs
src/HotChocolate/Core/src/Types/Execution/Processing/Tasks/BatchResolverTask.cs
\end{minted}

\index{resolver count!what chapter 3 was really counting}%
So the number chapter~\ref{ch:the-life-of-a-request} introduced counts resolver
tasks and not fields resolved, and it always did. Nothing was wrong with any
earlier chapter's arithmetic, because until this one every field on that path
did something asynchronous and the two numbers were the same. The moment a
resolver stopped needing to wait for anything, they parted company, and a metric
that had been reliable for ten chapters started measuring something adjacent to
what its name says.

That is the second time in this book that a number turned out to be answering a
slightly different question than the one it was asked. The first was
chapter~\ref{ch:schema-design-that-survives-change}'s discovery that \code{[ID]}
had been inert for three chapters. Both were found the same way: by changing
something and being surprised by which direction a number moved.

\medskip
Three of those numbers went down and none of them is a saving. The work went
somewhere, and where it went is the last thing this chapter has to account for.

\index{query plan!at six subgraphs|)}%
